\subsection{Geeignete Projekttypen}
\label{subsec:geeignete-projekttypen}

Die Eignung hängt von Zugriffskanälen, Backend und Persistenz, Deployment,
nativer Integration sowie Infrastruktur-, Sicherheits- und Betriebsbedarf ab.
Belegt sind generierbare Kombinationen und geprüfte Profilpfade. Die Stufen
\emph{gut geeignet}, \emph{bedingt geeignet} und \emph{nicht primär
vorgesehen} bewerten dagegen qualitativ den Abstand zur Baseline, nicht die
Erfolgswahrscheinlichkeit eines Projekts
\parencite{kleiveist2026profiles,kleiveist2026acceptance}.
Je mehr zentrale Anforderungen außerhalb der vorhandenen Komponenten und
Nachweise liegen, desto größer sind Anpassungsbedarf und verbleibendes Risiko.

\texttt{web-only} eignet sich für browserbasierte Frontends ohne
Template-Backend. \texttt{web-cloud} ergänzt zentrale API und providerneutrale
Deploymentstruktur, etwa für interne Business- und Datenanwendungen.
PostgreSQL bleibt optional; Authentifizierung, Berechtigungen, Fachmodelle und
produktive Infrastruktur sind produktspezifisch.
Damit deckt das Profil die technische Ausgangsstruktur, nicht die konkrete
Geschäftsanwendung ab.

\texttt{desktop-local} passt zu lokalen Werkzeugen ohne zentrales Backend,
sofern native Funktionen überschaubar ergänzt werden können; eine lokale
Datenbank fehlt. \texttt{desktop-cloud} ergänzt zentrale Dienste,
\texttt{full-platform} zusätzlich den Browserkanal. Der primäre Einsatzbereich
sind damit Web-, Cloud- und Desktop-orientierte Business-/Full-Stack-Anwendungen.
Für alle Desktopvarianten müssen weitergehende native Funktionen gezielt
ergänzt und geprüft werden.

\begin{figure}[H]
  \centering
  \resizebox{0.98\textwidth}{!}{%
  \begin{tikzpicture}[
    cell/.style={draw=caseLine,minimum height=11mm,align=left,anchor=center,
      font=\sffamily\scriptsize,text width=62mm,inner sep=3mm},
    head/.style={cell,fill=caseBlue,text=white,align=center,
      font=\sffamily\bfseries\scriptsize,text width=29mm},
    good/.style={head,fill=caseGreen,minimum height=21mm},
    conditional/.style={head,fill=caseBlue!82,minimum height=19mm},
    outside/.style={head,fill=black!62,minimum height=19mm},
    types/.style={cell,text width=57mm},
    reason/.style={cell,text width=67mm}
  ]
    \matrix[matrix of nodes,nodes in empty cells,ampersand replacement=\&,row sep=-\pgflinewidth,column sep=-\pgflinewidth] {
      |[head]| Kategorie \& |[head,text width=57mm]| Projekttypen \& |[head,text width=67mm]| Maßgebliches Kriterium \\
      |[good]| \shortstack{GUT\\GEEIGNET}
        \& |[types,minimum height=21mm]| Browser-Frontend; Web + Backend/Cloud; lokales Desktop-Werkzeug; Desktop + Cloud; Web + Desktop + Cloud
        \& |[reason,minimum height=21mm]| Zugriffskanäle und Laufzeitschichten entsprechen einem geprüften Profil; Fach- und Betriebsfunktionen bleiben zu ergänzen. \\
      |[conditional]| \shortstack{BEDINGT\\GEEIGNET}
        \& |[types,minimum height=19mm]| hoher nativer Anteil; stark individualisierte Infrastruktur; hohe Compliance- oder Sicherheitsauflagen
        \& |[reason,minimum height=19mm]| Ein Profil deckt nur einen Teil ab; erhebliche Erweiterungen oder zusätzliche Nachweise sind erforderlich. \\
      |[outside]| \shortstack{NICHT PRIMÄR\\VORGESEHEN}
        \& |[types,minimum height=19mm]| harte Echtzeitsysteme; Game-Engine-basierte Spiele; native Mobile-Apps; Firmware- und Embedded-Kernsysteme
        \& |[reason,minimum height=19mm]| Die tragende Laufzeit- oder Plattformarchitektur liegt außerhalb der Vite-/FastAPI-/Tauri-Baseline. \\
    };
  \end{tikzpicture}}
  \caption{Qualitative Eignungsübersicht nach Architekturfit und Anpassungsbedarf}
  \label{fig:project-suitability-map}
  \smallskip
  {\footnotesize Quelle: Eigene Transferbewertung auf Grundlage der
  Profildefinitionen und der technischen Abnahme
  \parencite{kleiveist2026profiles,kleiveist2026acceptance}.\par}
\end{figure}


Abbildung~\ref{fig:project-suitability-map} und
Tabelle~\ref{tab:project-suitability} fassen die Zuordnung zusammen.
\emph{Geeignet} bedeutet nur, dass die Baseline wesentliche Plattformteile
abdeckt. Fachlogik, Oberfläche, Integrationen, Produktsicherheit und Betrieb
bleiben ebenso offen wie die wirtschaftliche Bewertung.

\begin{table}[H]
  \centering
  \scriptsize
  \renewcommand{\arraystretch}{1.12}
  \caption{Zuordnung technischer Projekttypen zur Template-Baseline}
  \label{tab:project-suitability}
  \begin{tabularx}{\textwidth}{@{}>{\raggedright\arraybackslash}p{0.19\textwidth}>{\raggedright\arraybackslash}p{0.13\textwidth}>{\raggedright\arraybackslash}p{0.20\textwidth}>{\raggedright\arraybackslash}X@{}}
    \toprule
    \textbf{Projekttyp} & \textbf{Eignung} & \textbf{Profil} & \textbf{Notwendige Ergänzungen und Grenzen} \\
    \midrule
    Browser-Frontend & gut & \texttt{web-only}
      & Fachlogik und UI; externe APIs möglich, aber kein FastAPI-Backend der Baseline. \\
    Webanwendung mit API/Cloud & gut & \texttt{web-cloud}
      & Authentifizierung, Fachmodelle und produktiver Betrieb; PostgreSQL nur bei relationalem Persistenzbedarf. \\
    lokales Desktop-Werkzeug & gut & \texttt{desktop-local}
      & Native Kommandos, Packaging und Signing produktspezifisch; keine lokale Datenbank enthalten. \\
    Desktop mit zentralem Backend & gut & \texttt{desktop-cloud}
      & Produktive API, Security und Signing; PostgreSQL bleibt optional. \\
    Browser + Desktop + Cloud & gut & \texttt{full-platform}
      & Kanalbezogene UX und Tests; Persistenz nur mit separater Capability. \\
    Anwendung mit hohem nativen Anteil & bedingt & Desktopprofil oder andere Basis
      & Erheblicher Aufwand für Plattform-APIs, Berechtigungen, systemnahe Dienste oder native UI. \\
    stark spezialisierte Infrastruktur & bedingt & nach Teilarchitektur
      & Eignung nur, wenn Client oder Einzeldienst sinnvoll abgrenzbar ist; Infrastruktur nicht vorgegeben. \\
    Hard Real-Time, Game Engine, native Mobile- oder Embedded-Kernsysteme
      & nicht primär vorgesehen & kein Profil
      & Tragende Laufzeit- und Plattformanforderungen liegen außerhalb der Baseline. \\
    \bottomrule
  \end{tabularx}
  \smallskip
  {\footnotesize Quelle: Eigene qualitative Zuordnung auf Grundlage von
  \textcite{kleiveist2026profiles} und \textcite{kleiveist2026acceptance}.\par}
\end{table}

\beginnerinput{workspace/statement/800_einsatz_transfer/810}
